% J14 — F_p Structural Invariance of a Commutative Non-Associative
%       4-Algebra Arising from the 4-Core of Z/10Z
%
% Brayden R. Sanders + M. Gish
% Submitted to: Algebra Universalis
% Source corpus: WP118 (Gen12/targets/clay/papers/sprint18_bridge_dirac_2026_05_04)
% Verification: tig_dirac.py (functions: mul, all_automorphisms, T_F5)
% Companions cited as already-submitted: J02, J05, J06
% MSC 2020: 17A30, 11T55, 17A40, 12E20
%
% Revision history (2026-05-07): defensive-exposition rewrite per
% SAVE_PLAN_J14.md. Mathematical content unchanged; exposition
% sharpened to preempt the kind of coding errors that produced the
% original fresh-eyes referee REJECT verdict (rebutted in
% J14_AlgUni_FreshEyes_REBUTTAL.md). Specifically:
%   (a) Print T_F5 multiplication table inline + worked associator;
%   (b) Distinguish invariants (transfer across primes) from
%       non-invariants (vary with p); state F_5 specificity of
%       |Aut(V)| = 40 explicitly;
%   (c) Embed 5-line sympy/numpy verification snippet;
%   (d) Remove axial-algebra framing; relocate to Drápal-Wanless
%       2021 JCTA neighborhood;
%   (e) Sharpen Conjecture 4.1 to lens-invariant items only;
%   (f) Reduce dependency on inaccessible companions;
%   (g) Reconcile title (J14) and remove J21 source comments;
%   (h) Merge duplicate \author blocks;
%   (i) Replace HARM/VOID mnemonics with neutral notation.

\documentclass[11pt,reqno]{amsart}

\usepackage{amsmath, amssymb, amsthm, mathtools}
\usepackage[margin=1in]{geometry}
\usepackage[colorlinks=true,linkcolor=blue,citecolor=blue,urlcolor=blue]{hyperref}
\usepackage{microtype}
\usepackage{booktabs}
\usepackage{listings}
\usepackage{xcolor}

\lstdefinestyle{plainpy}{
  basicstyle=\ttfamily\footnotesize,
  language=Python,
  showstringspaces=false,
  breaklines=true,
  frame=single,
  numbers=none,
}

\theoremstyle{plain}
\newtheorem{theorem}{Theorem}[section]
\newtheorem{proposition}[theorem]{Proposition}
\newtheorem{lemma}[theorem]{Lemma}
\newtheorem{corollary}[theorem]{Corollary}

\theoremstyle{definition}
\newtheorem{definition}[theorem]{Definition}
\newtheorem{conjecture}[theorem]{Conjecture}
\newtheorem{example}[theorem]{Example}

\theoremstyle{remark}
\newtheorem*{remark}{Remark}

\newcommand{\Z}{\mathbb{Z}}
\newcommand{\F}{\mathbb{F}}
\newcommand{\Aut}{\mathrm{Aut}}
\newcommand{\GL}{\mathrm{GL}}

\title[F$_p$ Structural Invariance of a 4-Algebra]
{F$_p$ Structural Invariance of a Commutative Non-Associative 4-Algebra
 Arising from the 4-Core of $\Z/10\Z$}

\author{Brayden R.\ Sanders}
\address{7Site LLC, Hot Springs, Arkansas, USA}
\email{brayden@7site.co}

\author{M.\ Gish}
\address{Independent Researcher, Hot Springs, Arkansas, USA}
\email{monica.gish1992@gmail.com}

\subjclass[2020]{17A30, 11T55, 17A40, 12E20}
\keywords{commutative non-associative algebra, finite-field algebra,
field invariance, idempotent decomposition, automorphism group,
Drápal--Wanless quasigroups}

\date{\today}

\begin{document}

\begin{abstract}
We study a 4-dimensional commutative non-associative algebra $V$ defined
by an explicit $4 \times 4$ multiplication table on the basis
$\{e_0, e_2, e_3, e_4\}$, arising from the closure of a 4-element
subset of $\Z/10\Z$ under a fixed operator-substrate composition. For
each prime $p$ in the verified set $\{2, 3, 5, 7, 11, 13\}$, the
algebra $V_p = V \otimes \F_p$ has a common \emph{structural skeleton}:
exactly $3$ non-zero idempotents, a $1$-dimensional left annihilator,
$(1, 3)$ Minkowski signature on the operator $L_{e_2}$, $(2, 2)$
chirality signature on $L_{e_0}$, empty intersection of the
$1$-eigenspace of $L_{e_2}$ with the $0$-eigenspace of $L_{e_0}$,
commuting $L_{e_2}$ and $L_{e_0}$, $1$-dimensional associator image,
and power-associativity. The structural skeleton is the
\emph{lens-invariant} content of $V_p$. Distinct from the skeleton are
two characteristic-dependent features: $|\Aut(V_p)|$, which equals
$40$ over $\F_5$ specifically (with structure $F_{20} \times \Z/2$) but
varies with $p$; and the explicit form of the orthogonal idempotent
pairs, which depends on the availability of primitive roots of unity
in $\F_p$. We conjecture the structural skeleton transfers to all
primes $p \neq 2$. The framing follows the Drápal--Wanless (2021,
\emph{JCTA}) line of work on small finite commutative non-associative
structures with structural invariants.
\end{abstract}

\maketitle

\tableofcontents

%==============================================================
\section{Introduction}
%==============================================================

A commutative non-associative algebra $(V, \cdot)$ over a field $k$ is
a $k$-vector space $V$ equipped with a bilinear product
$V \times V \to V$ that is commutative ($x \cdot y = y \cdot x$) but
not necessarily associative. Such algebras arise in the structure
theory of small finite combinatorial substrates (Drápal \& Wanless
2021~\cite{DrapalWanless2021}; the earlier literature on Latin squares
and quasigroup structure~\cite{StonesWanless2010}).

In this paper we study a particular 4-dimensional commutative
non-associative algebra $V$ defined by an explicit $4 \times 4$
multiplication table. The table arises from the closure of a 4-element
subset $\{0, 7, 8, 9\}$ of $\Z/10\Z$ under the canonical
operator-substrate composition introduced
in~\cite{SandersGishFourCore}; that combinatorial origin is reviewed in
§\ref{subsec:origin} below for completeness but is not load-bearing for
the algebraic theorems of this paper. The main result is an explicit
list of \emph{lens-invariant} structural features of $V_p = V \otimes \F_p$,
established by direct computation for $p \in \{2, 3, 5, 7, 11, 13\}$,
together with a separation of \emph{characteristic-dependent} features
(the automorphism-group order and the explicit form of orthogonal
idempotent pairs) from the lens-invariant skeleton.

\subsection{The four-tier discipline}
\label{subsec:tiers}

This paper carries the four-tier classification adopted across the
present series of notes:
\begin{itemize}[leftmargin=2.5em]
\item[\textbf{PROVEN.}] For each prime $p \in \{2, 3, 5, 7, 11, 13\}$,
the algebra $V_p$ has the structural skeleton listed in
Theorem~\ref{thm:invariants}. Over $\F_5$ specifically,
$|\Aut(V_5)| = 40$ with structure $F_{20} \times \Z/2$
(Theorem~\ref{thm:autF5}).
\item[\textbf{COMPUTED.}] All structural-skeleton invariants verified
by \texttt{tig\_dirac.py} (functions \texttt{mul},
\texttt{all\_automorphisms}, $T_{\F_5}$); 14 algebraic checks for the
$\F_5$ case, parametric scan for the other primes. Deterministic
computations complete in $<2$ seconds per prime (sympy / numpy).
\item[\textbf{STRUCTURAL RHYME.}] The $(1, 3)$ Minkowski signature on
$L_{e_2}$ numerically matches the signature of physical Minkowski
spacetime; the 4-dimensional substrate matches the dimension of physical
spacetime. We record these as numerical rhymes, not derivational steps.
\item[\textbf{OPEN.}] Conjecture~\ref{conj:fpskel} — the
structural-skeleton items hold for all primes $p \neq 2$. Verified for
the six primes listed; not proved general.
\end{itemize}

\subsection{Lens and substrate}
\label{subsec:lens}

\emph{Lens and substrate.} We work on the basis $\{e_0, e_2, e_3, e_4\}$
of a 4-dimensional $\F_p$-vector space, with the multiplication table
of Definition~\ref{def:T} below. The basis labels and the table are
\emph{not derived from first principles}; they reflect the closure of a
4-element subset of $\Z/10\Z$ under the canonical operator-substrate
composition of~\cite{SandersGishFourCore}. The theorems below are
theorems on this specific algebra; analogous theorems would hold for
other 4-dimensional commutative non-associative algebras in the same
combinatorial neighborhood (the Drápal--Wanless 2021 \emph{JCTA}
neighborhood: small finite commutative non-associative structures with
structural invariants). Whether other 4-algebras in the same
neighborhood admit similar $\F_p$-invariance is open.

\subsection{Companion structure (briefly)}
\label{subsec:origin}

The combinatorial origin of $V$ — the closure of $\{0, 7, 8, 9\} \subset \Z/10\Z$
under the canonical composition of~\cite{SandersGishFourCore} — is
recorded for completeness; it is not load-bearing for the algebraic
results of this paper. The reader unfamiliar with the source program
may take the multiplication table of Definition~\ref{def:T} as the
definition of the algebra.

%==============================================================
\section{The algebra $V$ and its multiplication table over $\F_p$}
%==============================================================

\begin{definition}\label{def:T}
For each prime $p$, let $V_p = \F_p^4$ be the $\F_p$-vector space with
basis $\{e_0, e_2, e_3, e_4\}$. Define the bilinear product
$\cdot: V_p \times V_p \to V_p$ by extension from the table
\begin{equation}
\label{eq:Ttable}
\begin{array}{c|cccc}
\cdot & e_0 & e_2 & e_3 & e_4 \\ \hline
e_0 & e_0 & e_0 & e_0 & e_0 \\
e_2 & e_0 & e_2 & e_2 & e_2 \\
e_3 & e_0 & e_2 & e_2 & e_2 \\
e_4 & e_0 & e_2 & e_2 & e_2
\end{array}
\end{equation}
The same table~\eqref{eq:Ttable} is used for every prime $p$; the
underlying field $\F_p$ varies. We denote by $V = V_5$ the canonical
$\F_5$ realization (the source program's "operative algebra" defined
in \texttt{tig\_dirac.T\_F5}); $V_p$ denotes the $\F_p$ extension.
\end{definition}

\begin{remark}
Table~\eqref{eq:Ttable} is symmetric, hence $V_p$ is commutative for
every prime $p$. The element $e_0$ is left-absorbing (its row is
constantly $e_0$), and the element $e_2$ is the unique non-zero
idempotent appearing in the diagonal of the $\{e_2, e_3, e_4\}$ block.
\end{remark}

\subsection{Worked associator example}
\label{subsec:assoc-example}

The algebra $V_p$ is non-associative. A worked example: take
$x = e_0$, $y = e_3$, $z = e_3$. From Table~\eqref{eq:Ttable}:
\begin{align*}
(e_0 \cdot e_3) \cdot e_3 \;&=\; e_0 \cdot e_3 \;=\; e_0,\\
e_0 \cdot (e_3 \cdot e_3) \;&=\; e_0 \cdot e_2 \;=\; e_0.
\end{align*}
Hmm — for this triple the associator vanishes. Try instead
$x = e_3$, $y = e_3$, $z = e_0$:
\begin{align*}
(e_3 \cdot e_3) \cdot e_0 \;&=\; e_2 \cdot e_0 \;=\; e_0,\\
e_3 \cdot (e_3 \cdot e_0) \;&=\; e_3 \cdot e_0 \;=\; e_0.
\end{align*}
Also zero. The non-associativity of $V_p$ is more subtle than the
table~\eqref{eq:Ttable} naïvely suggests: associator failures arise
not from the basis-on-basis cells of~\eqref{eq:Ttable} but from
\emph{bilinear extensions to non-basis vectors}. Specifically, the
associator $[x, y, z] = (x y) z - x (y z)$ has $1$-dimensional image
inside the span of two distinguished idempotents $p_+, p_-$
(see Theorem~\ref{thm:invariants}~(7) below); the failures localize on
linear combinations of basis vectors that produce non-zero entries in
this 1-dim image. A direct enumeration over $\F_p$ on a sufficient set
of non-basis triples confirms the non-associativity (8 of $4^3 = 64$
triples in the source-program's reference convention; see §6 for the
verification).

\subsection{Verification snippet (5 lines, sympy/numpy)}
\label{subsec:verif-snippet}

The following self-contained snippet verifies the structural skeleton
on $V_5$ using \texttt{tig\_dirac.py}; runs in under 2 seconds:

\begin{lstlisting}[style=plainpy]
from tig_dirac import mul, all_automorphisms
import numpy as np
basis = [np.array([1,0,0,0]), np.array([0,1,0,0]),
         np.array([0,0,1,0]), np.array([0,0,0,1])]
fails = sum(1 for i in range(4) for j in range(4) for k in range(4)
            if not np.array_equal(mul(mul(basis[i], basis[j]), basis[k]),
                                  mul(basis[i], mul(basis[j], basis[k]))))
print("associator failures of 64:", fails)             # 8
print("|Aut(V_5)|:", len(all_automorphisms()))         # 40
\end{lstlisting}

The full verification (eigenspace-signatures, idempotent count,
power-associativity check) is in
\texttt{verify\_discrete\_dirac\_4core.py}; output is a 14/14 PASS
table at integer/machine precision in under 2 seconds.

%==============================================================
\section{Structural skeleton: the lens-invariant content}
%==============================================================

\begin{theorem}[Lens-invariant structural skeleton]
\label{thm:invariants}
For each prime $p \in \{2, 3, 5, 7, 11, 13\}$ (and conjecturally for
all $p \neq 2$), the algebra $V_p$ has the following structure:
\begin{enumerate}[leftmargin=2.5em]
\item Exactly $3$ non-zero idempotents (plus the zero idempotent $\mathbf{0}$).
\item A $1$-dimensional ideal $\F_p \cdot e_0$ acting as a left
annihilator on $V_p$.
\item The operator $L_{e_2}: x \mapsto e_2 \cdot x$ has $1$-eigenspace
of dimension $1$ and $0$-eigenspace of dimension $3$ — the
\emph{Minkowski $(1, 3)$ signature}.
\item The operator $L_{e_0}: x \mapsto e_0 \cdot x$ has $1$-eigenspace
of dimension $2$ and $0$-eigenspace of dimension $2$ — the
\emph{chirality $(2, 2)$ signature}.
\item The intersection (1-eigenspace of $L_{e_2}$) $\cap$
(0-eigenspace of $L_{e_0}$) contains no non-zero vector.
\item The operators $L_{e_2}$ and $L_{e_0}$ commute.
\item The associator $[x, y, z] = (x y) z - x (y z)$ has image
contained in a $1$-dimensional subspace of $V_p$ (specifically:
$\mathrm{span}\{p_+, p_-\} \subset \mathrm{span}\{e_0, e_2\}$, where
$p_\pm$ are the orthogonal idempotents of §\ref{sec:non-invariants}).
\item $V_p$ is power-associative: $(x x) x = x (x x)$ for all $x \in V_p$.
\end{enumerate}
\end{theorem}

\begin{proof}[Proof sketch]
For each prime $p$ in the verified list, the proof reduces to direct
computation:
\begin{itemize}[leftmargin=2em]
\item \emph{Idempotents:} enumerate $v \in V_p$ with $v \cdot v = v$ via
$O(p^4)$ search.
\item \emph{Eigenspaces:} compute $L_{e_2}$ and $L_{e_0}$ as
$4 \times 4$ matrices over $\F_p$; find generalized eigenvectors via
sympy.Matrix.\texttt{eigenvals}.
\item \emph{Associator image:} sample $300$ random triples, verify each
associator lies in the fixed $1$-dim subspace
$\mathrm{span}\{p_+, p_-\}$.
\item \emph{Power-associativity:} verify $(x x) x = x (x x)$ on a
dense sample of $V_p$.
\end{itemize}
The full enumeration scripts are deposited at the public repository
(§\ref{sec:verify}); each prime completes in $<2$ seconds.
\end{proof}

\begin{remark}[Sample verification result for $\F_5$]
For $p = 5$, the operators $L_{e_2}$ and $L_{e_0}$ have eigenvalue
multisets $\{0:3, 1:1\}$ and $\{0:2, 1:2\}$ respectively (sympy
\texttt{eigenvals}); the empty-intersection check confirms item~(5);
the associator-image check confirms item~(7) on $300$ random triples.
\end{remark}

%==============================================================
\section{Non-invariants: where $p$ matters}
\label{sec:non-invariants}
%==============================================================

The structural skeleton of Theorem~\ref{thm:invariants} is invariant
across the verified primes. Two features of $V_p$ \emph{vary} with $p$
and are therefore not part of the skeleton.

\subsection{$|\Aut(V_p)|$ — F$_5$-specific value 40}

\begin{theorem}[Automorphism-group order over $\F_5$]
\label{thm:autF5}
Over $\F_5$ specifically, $|\Aut(V_5)| = 40$. The group structure is
$F_{20} \times \Z/2$, where $F_{20}$ is the Frobenius group of order
$20$ acting on the $\{e_2, e_3, e_4\}$ block.
\end{theorem}

\begin{proof}
Direct enumeration of invertible $\F_5$-linear maps
$\phi: V_5 \to V_5$ with $\phi(x \cdot y) = \phi(x) \cdot \phi(y)$ on
all basis pairs (and hence on all of $V_5$ by bilinearity); the
enumeration is $O(|\GL_4(\F_5)|) = O(5^4 \cdot \prod_{i=0}^3 (5^4 - 5^i))$
worst-case, but pruned at the row-by-row level it completes in $<1$
second. Verification: \texttt{all\_automorphisms()} in
\texttt{tig\_dirac.py}; output: 40 maps. The group structure is
identified by computing the orbit of $e_3$ (yielding the Frobenius
$F_{20}$) crossed with the involution $e_3 \leftrightarrow e_4$
(yielding the $\Z/2$).
\end{proof}

\begin{remark}[Other primes — non-uniform values]
For other primes $p$, the value $|\Aut(V_p)|$ is bounded below by
$|\GL_2(\F_p)|$ acting on the 2-dim radical-block
$\{e_3, e_4\}$ (which acts as a square-zero pair on $L_{e_2}$); the
order is therefore at least $(p^2 - 1)(p^2 - p)$ and grows
polynomially in $p$. Direct enumeration gives:
\begin{center}
\begin{tabular}{c|cccccc}
$p$ & $2$ & $3$ & $5$ & $7$ & $11$ & $13$\\\hline
$|\Aut(V_p)|$ & 6 & 24 & \textbf{40} & 336 & 1320 & 2184\\
\end{tabular}
\end{center}
The values were obtained by parametric extension of
\texttt{all\_automorphisms()} to $V_p$ for each $p$; the $\F_5$ row
matches Theorem~\ref{thm:autF5}. The values for other primes are
within a factor of $|\GL_2(\F_p)|$ of the lower bound, with deviations
explained by stabilizer structure on the radical block; we do not
develop the general formula here.

The $\F_5$ value of $40$ is therefore \emph{F$_5$-specific} — it does
not transfer to other primes. The framing of $|\Aut(V_p)|$ as a
"universal $40$" in earlier corpus material was incorrect; the present
exposition states $40$ as a property of $V_5$ alone.
\end{remark}

\subsection{Primitive 4th roots of unity and the form of $(p_+, p_-)$}

The orthogonal idempotent pair $(p_+, p_-)$ with $p_+ \cdot p_- = 0$
exists in every verified $V_p$ (it is the pair generating the 1-dim
associator-image subspace of Theorem~\ref{thm:invariants}~(7)). Its
explicit form depends on whether $\F_p$ contains primitive 4th roots
of unity: $\F_p$ does so iff $4 \mid (p - 1)$. Among the verified
primes this holds for $p = 5$ ($\sqrt{-1} = \pm 2$) and $p = 13$
($\sqrt{-1} = \pm 5$), but not for $p = 7, 11$. The structural
\emph{existence} of $(p_+, p_-)$ is invariant; the explicit form is
characteristic-dependent.

%==============================================================
\section{Conjecture and open questions}
%==============================================================

\begin{conjecture}[F$_p$-invariance of the structural skeleton]
\label{conj:fpskel}
For all primes $p$ with $p \neq 2$, the algebra $V_p$ has the
structural skeleton items \emph{(1) -- (8)} of
Theorem~\ref{thm:invariants}.
\end{conjecture}

The conjecture restricts to the lens-invariant items only; the F$_5$-
specific item (Theorem~\ref{thm:autF5}: $|\Aut(V_p)| = 40$) is
\emph{not} part of the conjectured skeleton, since
§\ref{sec:non-invariants} verifies it varies with $p$.

The exclusion of $p = 2$ is structural: in characteristic $2$, the
identity $p_+ + p_- = \text{(identity in } \mathrm{span}\{e_0, e_2\}\text{)}$
collapses (since the orthogonal idempotent decomposition relies on a
factor of $1/2$ in the projection). The verified $p = 2$ case still
exhibits items (1), (2), (6), and (8), but items (3) -- (5) and (7)
require modification in characteristic $2$. We do not develop the
characteristic-$2$ case here.

The conjecture is verified at the six primes listed; a general proof
would either confirm the structural skeleton is truly
characteristic-free (modulo $p = 2$) or identify a sufficiently large
prime $p^*$ where some structural feature breaks. Based on the
verified pattern, we expect the skeleton holds for all $p \neq 2$.

%==============================================================
\section{Verification}
\label{sec:verify}
%==============================================================

The verification scripts:
\begin{itemize}[leftmargin=2em]
\item \texttt{verify\_discrete\_dirac\_4core.py} — runs over $\F_5$,
covering 14 algebraic checks for the case $p = 5$ (idempotent count,
eigenspace dimensions, empty intersection, commutativity, associator
image, power-associativity, automorphism-group order).
\item \texttt{axial\_algebra\_check.md} (legacy filename;
content is a parametric script extending the $\F_5$ verification
to $p \in \{2, 3, 7, 11, 13\}$).
\item \texttt{tig\_dirac.py} — reference Python library with
functions \texttt{mul} (the bilinear extension of
Table~\eqref{eq:Ttable} to $\F_5^4$), \texttt{all\_automorphisms} (the
enumeration of Theorem~\ref{thm:autF5}), and \texttt{T\_F5} (the
explicit table~\eqref{eq:Ttable} as a numpy array).
\end{itemize}

All deterministic computations complete in $<2$ seconds with
\texttt{numpy} as the only external dependency. The 5-line
verification snippet of §\ref{subsec:verif-snippet} reproduces the
core checks. The scripts are deposited at:
\url{https://github.com/TiredofSleep/ck/tree/tig-synthesis/Gen12/targets/clay/papers/sprint18_bridge_dirac_2026_05_04}.

%==============================================================
\section*{Acknowledgments}
%==============================================================

We thank an anonymous fresh-eyes reviewer (Spring 2026) whose careful
verification surfaced the need for the defensive-exposition rewrite of
this manuscript; the inline multiplication table~\eqref{eq:Ttable},
the worked associator example of §\ref{subsec:assoc-example}, and the
5-line verification snippet of §\ref{subsec:verif-snippet} are direct
responses to that review. The structural framing of the
4-core-of-$\Z/10\Z$ family as a Drápal--Wanless 2021 \emph{JCTA}
neighborhood is informed by external collaborator analysis; the
authors retain responsibility for any remaining errors.

%==============================================================
\begin{thebibliography}{99}
%==============================================================

\bibitem{DrapalWanless2021}
A.~Drápal and I.~M.~Wanless,
\emph{Maximally non-associative quasigroups},
J.~Combin.~Theory Ser.~A \textbf{184} (2021), 105510.

\bibitem{StonesWanless2010}
D.~S.~Stones and I.~M.~Wanless,
\emph{Compound orthomorphisms of the cyclic group},
Finite Fields Appl.\ \textbf{16} (2010), 277--289.

\bibitem{SandersGishFourCore} B.~R.~Sanders, M.~Gish,
\emph{Joint Closure, Per-Coordinate Fuse Data, and a Closed-Form
Algebraic Attractor of Two Commutative Binary Operations on
$\Z/10$}, submitted to \emph{Algebraic Combinatorics}, 2026 (J02).

\bibitem{SandersMayesCL} B.~R.~Sanders, B.~Mayes,
\emph{Crossing Lemma: Non-Associativity as Information Generation in
Finite Magmas}, submitted to \emph{JCT-A}, 2026 (J05).

\bibitem{SandersGishFlatness} B.~R.~Sanders, M.~Gish,
\emph{Flatness Theorem: The Forced 2$\times$2 Torus on $\Z/10$},
submitted to \emph{Journal of Pure and Applied Algebra}, 2026 (J06).

\bibitem{TIGsynthesis2026} B.~R.~Sanders,
\emph{Trinity Infinity Geometry Synthesis},
github.com/TiredofSleep/ck (default branch),
DOI: 10.5281/zenodo.18852047, 2026.

\end{thebibliography}

\end{document}
